\scanpage{50}
\begin{flushright}
at\illegible{1} \mspage{10}
\end{flushright}

\noindent fixed values of $\alpha$,and that even these limited results depend on computation beyond any reasonable prospect of hand calculation(except for $\alpha=1$). However in view of the complete test of the M-results the prggrams can be regarded as self-checking,since the \unsure{M|R|A}\margin{\handwritten{\tikz[x=1em,y=1em,baseline=-.1em]{\draw (.5,.72)--(.08,-.35) (.27,-.10)--(.38,-.35);} $-$\\(hyphen)}}\deleted{\unsure{;|j|,}}\rlap{\smash{\handwritten{\tikz[x=1em,y=1em,baseline=-.1em]{\draw (.10,-.45)--(.60,.95) (.29,-.22)--(.40,-.45);}}}}results follow theoretically from the independently computed values of $\handwritten{c_{ij}}$ in \handwritten{(4.1)} .

Since the subgroups $\handwritten{\Gamma_{0}}(13)$, $\handwritten{\Gamma_{0}}(17)$, $\handwritten{\Gamma_{0}}(19)$,and $\handwritten{\Gamma_{0}}(23)$ of the modular group have genera $0,1,1,\text{and }2$ respectively,it is remarkable that all these primes should exhibit the same behaviour with regard to our multiplicative congrue\makebox[0pt][l]{\handwritten{\tikz[x=1em,y=1ex,baseline=0pt]{\draw (.02,-.28)--(.54,-.28)--(.54,1.54)--(1.01,1.6);}}}cne p\makebox[0pt][l]{\handwritten{\tikz[x=1em,y=1ex,baseline=0pt]{\draw (.02,-.28)--(.54,-.28)--(.54,1.54)--(1.01,1.6);}}}orperties. It thus seems possible that some entirely different method or theory may exist which would give an uniform proof of all these cases.

\deleted{\unsure{4|3}}\handwritten{5}. \uline{The} \uline{cases} $q\handwritten{\geqslant}29$. As the reader may have anticipated from the non-appearance of $29$ in \handwritten{\S\,4},the direct analogy to \handwritten{(3.9)} and \handwritten{(3\unsure{.|:}10)} breaks down for $q=29$. Considering first the R-results modulo $29$ only,we can prove by the combination of modular theory and computing described in \handwritten{\S}\deleted{\unsure{3|5|4}}\handwritten{4} that
\[
\handwritten{\begin{aligned}
U_{29}\,j(x)&\equiv a_{1}\,g_{1}(x)+a_{2}\,g_{2}(x)\quad(\mathrm{mod}\;29),\\
U_{29}\,g_{1}(x)&\equiv 26\,g_{1}(x)+14\,g_{2}(x)\quad(\mathrm{mod}\;29),\\
U_{29}\,g_{2}(x)&\equiv 3\,g_{1}(x)+21\,g_{2}(x)\quad(\mathrm{mod}\;29).
\end{aligned}}
\]
Thus
\[
\resizebox{\linewidth}{!}{$\displaystyle \handwritten{U_{29}\{g_{1}(x)+\mu\,g_{2}(x)\}\equiv(3\mu-3)\,g_{1}(x)+(14-8\mu)\,g_{2}(x)\quad(\mathrm{mod}\;29),}$}
\]
and if we attempt to make the RHS a multiple of $\handwritten{\{}g_{\handwritten{1}}(x)\ \deleted{\text{and }g_{2}}\handwritten{{}+\mu\,g_{2}(x)\}}$ we are led to the congruence
\[
\handwritten{\mu^{2}-8\mu+5\equiv 0\;(\mathrm{mod}\;29)}
\]
which has in fact no roots. However if we now consider our various coefficients\reversemarginpar\margin{\inserted{\handwritten{not}} \figure{\tikz[x=1em,y=1em,baseline=-.1em]{\draw (.10,-.45)--(.55,.95); \draw (.17,-.20) .. controls (.35,.02) and (.30,-.42) .. (.58,-.40);}}}\rlap{\smash{\handwritten{\tikz[x=1em,y=1em,baseline=-.1em]{\draw (.10,-.45)--(.40,.95); \draw (.15,-.18) .. controls (.31,.02) and (.34,-.44) .. (.58,-.42);}}}} as integers modulo $29$,but as elements of the Galois field $\mathrm{GF}(29^{2})$ we may assert(using ``$=$'' rather than ``$\handwritten{\equiv}$'' in a field of characteristic $29$):
\normalmarginpar
% archive sheet number: 49
